\documentclass[]{exam}
\usepackage{amssymb}
\usepackage{amsfonts}
\usepackage{amsmath}
\usepackage{amsthm}
\usepackage[top=1in, bottom=1.5in, left=1in, right=1in]{geometry}
\usepackage{tikz}
\usepackage{pgfplots}
\usepackage{empheq}
\usepackage{mathrsfs}
\checkboxchar{$\Box$}
\usepackage{graphicx}
\usepackage{caption}
\usepackage{float}
\usepackage{enumerate}
\usepackage{comment}
\usepackage{multicol}
\usepackage{color}
\usepackage{advdate}
\usepackage{gensymb}
\usetikzlibrary{arrows}
\usepackage{multirow}
\usepackage{arydshln}
\usepackage{verbatim}

\usepackage{pgfkeys,pgfcalendar}

\newcount\pgfdatecount
\newcommand{\tomorrow}{%
\pgfcalendardatetojulian{\year-\month-\day+1}{\pgfdatecount}
\pgfcalendarjuliantodate{\the\pgfdatecount}{\myyear}{\mymonth}{\myday}
\pgfcalendarmonthname{\mymonth}\space\myday,\space\myyear%
}

\newcommand{\advanceday}[1][1]{
\begingroup
\AdvanceDate[#1]
\today
\endgroup}

\newcommand{\vertex}{\node[vertex]}
\tikzstyle{vertex}=[circle, draw, inner sep=0pt, minimum size=6pt]
\usepackage{lastpage} 
\usepackage{extramarks}
\usepackage{graphicx} 
\usepackage{tikz}
\usepackage{sectsty} 
\usepackage{amsmath,amsfonts,amsthm} 
\usepackage{wrapfig, blindtext}

% Margins
\topmargin=-0.45in
\evensidemargin=0in
\oddsidemargin=0in
\textwidth=6.5in
\textheight=9.0in
\headsep=0.25in 
\numberwithin{equation}{section}
\linespread{1.1} 
\setlength\parindent{0pt} 
\setlength\answerskip{-.075in}
\setlength\answerlinelength{.5in}
\newcommand{\myitem}{\item[-]}

% HEADER AND FOOTER
\pagestyle{headandfoot}
\runningheadrule
\runningfootrule
\firstpageheader{Kutztown University}{  }{MATH 210}
\firstpageheadrule
\runningheader{P3.HW}{ }{MATH 210}
\firstpagefooter{}{}{} 
\runningfooter{}{}{Page \thepage\ of \numpages}


% TITLE PAGE
\newcommand{\assignmentname}{Phase 3 -- Homework} % Assignment title
\newcommand{\coursetitle}{MATH 210 -- Mathematical Computing and Typesetting}
\newcommand{\duedate}{Due: Friday, April 10 by 11:59 PM on D2L} % Update as needed
\newcommand{\readsections}{\, } % Teacher/lecturer
\newcommand{\semester}{\, }
\title{
\vspace{2in}
\textmd{\textbf{\course:\ \assignmentname}}\\
\normalsize\vspace{0.1in}{\duedate}\\
\vspace{0.1in}\large{\textit{\instructor}}
\vspace{3in}
}
\author{\textbf{Author: \authorname}}
\date{Last Updated: \today}

\pointsinrightmargin
\bracketedpoints



\begin{document}
%%%%%%%%%%%%%%%%%%%%%%%%%%%%%%%%%%%%%%%%%%%%%%%%%%%%%
% COMPILE TITLES AND SUCH
%%%%%%%%%%%%%%%%%%%%%%%%%%%%%%%%%%%%%%%%%%%%%%%%%%%%%
\hspace{1in}\\ \vspace{.15in}


\begin{center}
\normalsize
 \fbox{\fbox{\parbox{5.5in}{
        {\bf Instructions:}  Upload the \texttt{tex} file for this document to your Overleaf project.  All solutions shall be completed on this packet by typing your solution in \LaTeX\, in the space provided.  All solutions should be detailed and should clearly demonstrate the process by which you arrived at the answer.  Submit the compiled \texttt{pdf} file to D2L by 11:59 PM on the due date below.  Submit only a single pdf file of your entire packet.  The question will also ask you to make calculations in Python.  Upload any Python files to the appropriate folder in your GitHub repository.  In this folder, a \texttt{py} file is to be submitted for each problem such that when the \texttt{py} file is executed, the output (as presented in Python) is the solution to the problem. Academic dishonesty will not be tolerated.  }}}
\end{center}

\vspace{.5cm}

\begin{center}
	\Huge\textsc{
	\assignmentname}\\[4pt]
        \Large{\textsc{\coursetitle}\\[10pt]}
        \Large{\textsc{\duedate}\\[8pt]}
        \Large{\textsc{\readsections}\\[8pt]}
\end{center}

\vspace{1cm}

%%%%%%%%%%%%%%%%%%%%%%%%%%%%%%%%%%%%%%
% Uncomment the syntax below when you type up the solutions and type 
% your name.
\begin{center}
	\large{\textsc{\textcolor{red}{Solutions by [YOUR NAME HERE]}}}
\end{center}
\vfill

%%%%%%%%%%%%%%%%%%%%%%%%%%%%%%%%%%%%%%
% Reminders:
%%%%%%%%%%%%%%%%%%%%%%%%%%%%%%%%%%%%%%
Remember to begin each \texttt{py} file with \texttt{import numpy as np} and \verb|import matplotlib.pyplot as plt|.  
For this assignment you will also need \texttt{import networkx as nx}.  
If you compute eigenvalues/eigenvectors in Python, you may use \texttt{np.linalg.eig}.  

% Blank Page
\clearpage

\begin{questions}
%%%%%%%%%%%%%%%%%%%%%%%%%%%%%%%%%%%%%%
% NUMBER 1 
%%%%%%%%%%%%%%%%%%%%%%%%%%%%%%%%%%%%%%
\question Consider the matrix
\[
A=\begin{bmatrix}
5 & 2\\
1 & 2
\end{bmatrix}.
\]

\begin{parts}
\part Compute the characteristic polynomial $\det(A-\lambda I)$ and find the eigenvalues of $A$.  Type your work explicitly below.

\vfill

\part For each eigenvalue, find a corresponding eigenvector.  Type your work explicitly below and clearly label which eigenvector corresponds to which eigenvalue.

\vfill

\part Verify your answers by checking that $A\mathbf{v}=\lambda \mathbf{v}$ for each eigenpair $(\lambda,\mathbf{v})$.

\vfill

\part In a Python file, define $A$ as a NumPy array.  Use \texttt{np.linalg.eig} to compute eigenvalues and eigenvectors.  Print the eigenvalues/eigenvectors.

\end{parts}

\pagebreak

%%%%%%%%%%%%%%%%%%%%%%%%%%%%%%%%%%%%%%
% NUMBER 2
%%%%%%%%%%%%%%%%%%%%%%%%%%%%%%%%%%%%%%
\question Consider the following three matrices.  

Let the nodes be
\[
\{A,B,C,D,E\}.
\]
You are given three symmetric weighted adjacency matrices:
\[
A_{\text{strong}}=
\begin{bmatrix}
0 & 4 & 2 & 0 & 0\\
4 & 0 & 1 & 2 & 0\\
2 & 1 & 0 & 3 & 0\\
0 & 2 & 3 & 0 & 2\\
0 & 0 & 0 & 2 & 0
\end{bmatrix},
\qquad
A_{\text{medium}}=
\begin{bmatrix}
0 & 0 & 0 & 2 & 0\\
0 & 0 & 1 & 1 & 3\\
0 & 1 & 0 & 0 & 1\\
2 & 1 & 0 & 0 & 0\\
0 & 3 & 1 & 0 & 0
\end{bmatrix},
\qquad
A_{\text{weak}}=
\begin{bmatrix}
0 & 0 & 1 & 0 & 1\\
0 & 0 & 0 & 0 & 0\\
1 & 0 & 0 & 0 & 2\\
0 & 0 & 0 & 0 & 1\\
1 & 0 & 2 & 1 & 0
\end{bmatrix}.
\]

Define the combined interaction matrix
\[
A(\alpha)=A_{\text{strong}}+\alpha A_{\text{medium}}+\alpha^2 A_{\text{weak}}.
\]

\begin{parts}

\part Compute the strength of \emph{every} node, $s_A(\alpha),s_B(\alpha),s_C(\alpha),s_D(\alpha),s_E(\alpha)$, and simplify each expression.  Type them below.

\vfill

\part Find the intersection parameter $\alpha$ such that the strengths of nodes $C$ and $E$ are equal, i.e.,
\[ s_C(\alpha)=s_E(\alpha). \]
Solve exactly for $\alpha$ and clearly show your algebra below.  (Note: $\alpha$ could be bigger than 1.)

\vfill

\part Follow the code in the Python file called \verb|bemerick_p3_ws2_1.py|.  Create a new file that successfully performs the following:
\begin{enumerate}[(i)]
\item Define $A_{\text{strong}}$, $A_{\text{medium}}$, and $A_{\text{weak}}$ as \verb|numpy| arrays.
\item Define functions that return the strength of any two nodes, say $u$ and $v$.
\item Plot the two quadratic functions.
\item Confirm numerically the value of $\alpha$ you found in part (c) by using \verb|root_scalar| to find the intersection of $s_C(\alpha)$ and $s_E(\alpha)$.  
\item Modify the function \verb|plot_weighted_graph| to plot the weighted graph $A(\alpha)$ by introducing another input $A_{medium}$.  Plot the graph of $A(0.5)$ and include it in the \LaTeX\,report below.
\end{enumerate}

\end{parts}

\pagebreak

%%%%%%%%%%%%%%%%%%%%%%%%%%%%%%%%%%%%%%
% NUMBER 3
%%%%%%%%%%%%%%%%%%%%%%%%%%%%%%%%%%%%%%
\question Consider the undirected graph $G$ with vertex set
\[ V=\{A,B,C,D,E,F\} \]
and edge set
\[ E=\{AB,AC,BC,BD,CE,DE,DF,EF\}. \]

\begin{parts}
\part Use Python to plot this graph and include it below.  
\vfill
\part In the same Python file, compute the closeness, betweenness, and eigenvector centrality.  Include them all below and comment on if they agree.  Which node is the ``most important?" 
\vfill

\end{parts}

\end{questions}
\end{document}